% =====================================================================
%  CAPÍTULO 1 — INTRODUÇÃO
% =====================================================================
\chapter{Introdução}
\label{cap:introducao}

\section{Contextualização}
\label{sec:contextualizacao}

O interesse global pela genealogia e pela história familiar cresceu de
forma acentuada nas últimas décadas. A digitalização massiva de registos
históricos e a democratização do acesso à internet permitiram que
praticamente qualquer pessoa pudesse investigar as suas origens; em
paralelo, os testes de \textsc{dna} de ascendência abriram caminho para
que muitos, com origens até então desconhecidas, encontrassem finalmente
respostas. Plataformas como a \emph{Ancestry}, a \emph{MyHeritage}, a
\emph{FamilySearch} e a \emph{Findmypast} agregam hoje milhares de milhões
de registos e sustentam uma verdadeira era moderna de descoberta
familiar~\cite{familytreemag}.

Paralelamente a este movimento, a esmagadora maioria das memórias
familiares contemporâneas existe já em formato digital — fotografias no
telemóvel, vídeos no computador, documentos digitalizados. Contudo, esse
material encontra-se tipicamente \emph{disperso}, \emph{desorganizado} e
\emph{vulnerável}: vive espalhado por dispositivos, serviços de nuvem e
álbuns esquecidos, sem datação coerente, sem identificação das pessoas e,
sobretudo, sem qualquer fio narrativo que lhe dê sentido. À medida que os
membros mais velhos de uma família desaparecem, perde-se com eles o
contexto que tornava aquelas imagens significativas.

Os avanços recentes em \textbf{inteligência artificial generativa} — em
particular os \emph{Large Language Models} (\textsc{llm}) e os modelos
multimodais capazes de interpretar imagens — criam, pela primeira vez, a
possibilidade técnica de transformar este arquivo passivo num produto
narrativo vivo. É neste cruzamento entre preservação da memória familiar
e IA generativa que se enquadra o presente projeto.

\section{Definição do Problema}
\label{sec:problema}

As soluções existentes no mercado abordam apenas parcialmente o
problema. As grandes plataformas de genealogia são excelentes a
\emph{agregar registos} e a \emph{construir árvores}, mas não produzem
narrativa. Ferramentas de consumo como o \emph{Google~Photos} ou o
\emph{Apple~Memories} geram \emph{slideshows} automáticos visualmente
agradáveis, mas limitam-se a selecionar fotografias por critérios
algorítmicos, sem construir uma história coerente nem integrar dados de
múltiplas fontes (fotografias, documentos, árvores genealógicas). O que
falta — e que constitui o problema central deste trabalho — é a
\textbf{profundidade narrativa}: a capacidade de cruzar material
heterogéneo e dele extrair, de forma automática, uma história
factualmente ancorada, emocionalmente significativa e linguisticamente
correta.

Formalmente, o problema pode enunciar-se assim: \emph{dado um conjunto
heterogéneo e não estruturado de material de arquivo familiar
(fotografias com metadados, documentos digitalizados e uma árvore
genealógica), como produzir automaticamente — sem intervenção manual de
escrita — uma narrativa textual coerente e um vídeo documental narrado,
factualmente fiéis ao material fornecido e expressos em português
europeu?}

Este problema levanta vários desafios técnicos não triviais, que se
discutem em detalhe no \Cref{cap:estado_arte}: a \emph{alucinação} dos
modelos generativos (invenção de factos), a \emph{coerência temporal}
entre fontes com datas conflituosas, a \emph{identificação de pessoas} no
material e a própria \emph{avaliação da qualidade} narrativa.

\section{Objetivos}
\label{sec:objetivos}

\subsection{Objetivo geral}

Conceber, desenvolver e validar um sistema que receba material familiar
já existente e o transforme automaticamente numa narrativa multimédia —
concretamente, num vídeo documental narrado — recorrendo a inteligência
artificial generativa, com uma interface acessível e intuitiva.

\subsection{Objetivos específicos}

\begin{enumerate}[label=\textbf{O\arabic*.},leftmargin=*]
  \item \textbf{Ingerir e compreender material multimodal}: processar
  fotografias (metadados \textsc{exif} e conteúdo visual), documentos
  digitalizados (\textsc{ocr}) e árvores genealógicas (\textsc{gedcom}).
  \item \textbf{Estruturar temporalmente} os factos extraídos, resolvendo
  datas incompletas ou conflituosas e modelando explicitamente as
  relações de parentesco num grafo familiar.
  \item \textbf{Gerar narrativas} fluidas e factualmente ancoradas, em
  português europeu, controlando o tom e a estrutura através de
  \emph{templates} temáticos e de técnicas de \emph{prompt engineering}.
  \item \textbf{Produzir vídeos documentais} que combinem as fotografias
  enquadradas (com movimento opcional, efeito \emph{Ken~Burns}), a narração
  por síntese de voz e elementos cinematográficos (títulos, legendas,
  transições).
  \item \textbf{Disponibilizar} estas capacidades através de uma aplicação
  Web moderna, segura, multilingue e implantada na nuvem.
  \item \textbf{Garantir baixo custo e privacidade opcional}, assentando em
  ferramentas de código aberto e em camadas gratuitas, e mantendo a
  inferência \emph{local} como alternativa para operação privada e
  \emph{offline}.
\end{enumerate}

\section{Resultados Desejados e Âmbito}
\label{sec:ambito}

O âmbito do projeto, definido em articulação com a orientação, fixa
quatro tipos de \emph{input} — fotografias, vídeos, documentos e
informação genealógica — e um \emph{output} principal: um pequeno vídeo
documental ou \emph{slideshow} narrado, gerado automaticamente a partir
desse material. Dadas as limitações tecnológicas (e de custo) da geração
de vídeo longo por IA, o projeto concentra-se deliberadamente em
\emph{incrementos narrativos curtos} — segmentos de história que podem,
eventualmente, ser combinados numa narrativa mais longa. Esta opção é
academicamente defensável: um \emph{slideshow} narrado de alta qualidade,
construído programaticamente, cumpre integralmente os objetivos sem
depender de APIs de geração de vídeo dispendiosas.

\section{Contribuições}
\label{sec:contribuicoes}

As principais contribuições deste trabalho são:

\begin{itemize}[leftmargin=*]
  \item Uma \textbf{arquitetura modular em \emph{pipeline}} (M1--M4) que
  separa claramente ingestão, organização temporal, geração textual e
  produção multimédia.
  \item Um \textbf{subsistema RAG multi-inquilino com degradação
  graciosa}, que mantém o sistema funcional mesmo sem base vetorial.
  \item Um conjunto de \textbf{\emph{templates} de geração calibrados para
  português europeu}, que mitigam a tendência dos modelos para o
  português do Brasil.
  \item Uma \textbf{estratégia de IA em cascata e a custo nulo} (Groq na
  nuvem, Ollama local e Gemini), que concilia desempenho, robustez por
  \emph{fallback} e privacidade opcional (inferência local).
  \item Uma \textbf{narrativa segmentada em cenas} que liga cada parágrafo
  às fotografias que o ilustram, permitindo sincronizar imagem e narração
  no vídeo — a abordagem que distingue um documentário de um simples
  \emph{slideshow}.
  \item Uma \textbf{aplicação Web completa e implantada}, do
  \emph{frontend} ao \emph{deployment} na nuvem, com \textbf{geração
  assíncrona} (Celery ou executor \emph{in-process}) e resiliência a
  \emph{cold starts}.
\end{itemize}

\section{Estrutura do Documento}
\label{sec:estrutura}

O presente relatório encontra-se organizado em oito capítulos. O
\textbf{\Cref{cap:introducao}} (este) apresenta o contexto, o problema e
os objetivos. O \textbf{\Cref{cap:estado_arte}} analisa as soluções
existentes e enquadra teoricamente os conceitos fundamentais (LLM,
multimodalidade, RAG, geração de narrativa e de vídeo). O
\textbf{\Cref{cap:metodologia}} descreve a metodologia de desenvolvimento
e justifica a \emph{stack} tecnológica. O \textbf{\Cref{cap:requisitos}}
formaliza os requisitos funcionais e não funcionais e os casos de uso. O
\textbf{\Cref{cap:arquitetura}} apresenta a arquitetura do sistema, o
modelo de dados e o \emph{design} de interface. O
\textbf{\Cref{cap:implementacao}} detalha a implementação dos quatro
módulos e as decisões técnicas, com excertos de código críticos. O
\textbf{\Cref{cap:testes}} cobre a validação e os testes. Por fim, o
\textbf{\Cref{cap:conclusoes}} sintetiza os resultados, discute as
limitações e aponta o trabalho futuro.
